% ch03 译文：原文 §3a–3b（main.tex 第 516–810 行）
\section{代数准备\label{sec:algebraic-preparations}}
在铺开细节之前，本节先回顾 $A_{\infty}$-范畴的一些基本定义。全文除另行说明外，$\field$ 均表示特征为 $0$ 的域。
\subsection{$A_{\infty}$-范畴与 dg 范畴}
我们先回顾 $A_{\infty}$-范畴以及它们之间函子的定义。
\begin{dfn}
一个\emph{$\field$ 上的 $\mathbb{Z}$-分次 $A_{\infty}$-范畴 $\mathcal{A}$} 是指以 $\field$-向量空间的上链复形为 Hom、在 $A_{\infty}$-operad 上丰富的范畴。具体地说，它由以下数据构成：
\begin{itemize}
\item
\emph{对象}构成一个类 $\Ob(\mathcal{A})$。
% [原文 p.15]
\item 对每一对对象 $L,L'\in\Ob (\mathcal{A})$，一个 $\mathbb{Z}$-分次的 $\field$-向量空间，即态射空间
$$
\hom_{\mathcal{A}}(L,L')\coloneqq \bigoplus_{\bullet\in\mathbb{Z}} \hom_{\mathcal{A}}^{\bullet}(L,L').
$$
\item 一个\emph{$A_{\infty}$-结构} \[
    \mu_{\mathcal{A}}^{\bullet} =\left(\mu_{\mathcal{A}}^k\colon\hom_{\mathcal{A}}^{\bullet} (L_{k-1},L_k)\otimes\dotsb\otimes\hom_{\mathcal{A}}^{\bullet} (L_0, L_1)\to\hom_{\mathcal{A}}^{\bullet +2-k} (L_0, L_k)\right)_{k\in\mathbb{Z}_{\geq 0}}
\] 满足\emph{Maurer--Cartan 方程} $(\mu_{\mathcal{A}}^{\bullet})^2=0$，即 \[
    \sum_{m,n}(-1)^{\maltese_n}\mu_{\mathcal{A}}^{d-m+1} (a_d,\dotsb, a_{n-m+1},\mu_{\mathcal{A}}^m (a_{n+m},\dotsb, a_{n+1}),a_n,\dotsb, a_1)=0,
\] 其中 $\maltese_n =\vert a_1\vert+\dotsb+\vert a_n\vert -n$。特别地，$\mu_{\mathcal{A}}^1$ 是\emph{微分}，即对任意 $f\in\hom_{\mathcal{A}}(L,L')$ 有 $\mu_{\mathcal{A}}^1(\mu_{\mathcal{A}}^1(f))=0$。\[
\mu_{\mathcal{A}}^2\colon\hom_{\mathcal{A}}^{\bullet} (L_1,L_2)\otimes\hom_{\mathcal{A}}^{\bullet} (L_0,L_1)\to\hom_{\mathcal{A}}^{\bullet} (L_0, L_2)
\] 充当两个态射的\emph{复合}；此时 Maurer--Cartan 方程化为 Leibniz 法则 $\mu_{\mathcal{A}}^2 (\mu_{\mathcal{A}}^1 (f),g)\pm\mu_{\mathcal{A}}^2 (f,\mu_{\mathcal{A}}^1 (g))\pm\mu_{\mathcal{A}}^1\mu_{\mathcal{A}}^2 (f,g)=0$ 与同伦结合律 \begin{equation}\label{eq:homotopy-associativity}
\mu^1_{\mathcal{A}}\mu_{\mathcal{A}}^3 (f,g,h)=\pm\mu_{\mathcal{A}}^2 (\mu_{\mathcal{A}}^2 (f,g),h)\pm\mu_{\mathcal{A}}^2 (f,\mu_{\mathcal{A}}^2 (g,h)).
\end{equation}
为简便起见，对任意对象 $L_0$、$L_1$、$L_2\in\Ob(\mathcal{A})$，我们把 $(-1)^{-\vert f\vert}\mu_{\mathcal{A}}^2 (f',f)$ 也写作
$$
\comp
:
\hom_{\mathcal{A}}^{\bullet}(L_0,L_1)\underset{\field}{\otimes}\hom_{\mathcal{A}}^{\bullet}(L_1,L_2)
\rightarrow
\hom_{\mathcal{A}}^{\bullet}(L_0,L_2),
\quad
f\otimes f'
\mapsto
f\comp f'
$$
\end{itemize}
当 $\mu_{\mathcal{A}}^k=0$（$k\geq 2$）时，称 $\mathcal{A}$ 为\emph{dg 范畴}；此时同伦结合律 \eqref{eq:homotopy-associativity} 退化为严格结合律，于是 $\mathcal{A}$ 是真正的 $1$-范畴。在这种情形下，我们有时直接记 $d_{\mathcal{A}}\coloneqq\mu_{\mathcal{A}}^1$。
\end{dfn}

\begin{yizhu}[原文拼写与笔误]
本节原文有几处笔误，译文按原意处理，读者按下列对应理解即可：开头「recollment」应为「recollection」（回顾）；Leibniz 法则被拼作「Liebnitz」；dg 范畴的刻画写作 $\mu_{\mathcal{A}}^k=0$（$k\geq 2$），应为 $k\geq 3$——否则连复合 $\mu^2$ 也会被消掉，而 §3c 原文自己写的就是 $d\geq 3$。这些笔误不影响数学内容。
\end{yizhu}

\begin{conceptbox}{dg 范畴与 $A_{\infty}$ 范畴}
dg 范畴就是 Hom 复形带结合乘法（即 $\mu^2$）且 $\mu^k=0$（$k\geq 3$）的 $A_{\infty}$-范畴；$A_{\infty}$-范畴则允许复合只在同伦意义下结合，$\mu^3,\mu^4,\dots$ 逐阶记录结合律的失效。二者的差别在链层面，取上同调后都得到普通范畴 $H^{\bullet}(\mathcal{A})$。Fukaya 范畴天然是 $A_{\infty}$ 范畴（盘图展开给出高阶运算 $\mu^k$），而凝聚层复形天然是 dg 范畴；本文需要在同一层面上陈述函子与拟等价，故一律使用 $A_{\infty}$/dg 语言，并通过「高阶运算为零」在两者之间互相转换。
\par\smallskip\noindent{\footnotesize 参考：nLab，\emph{dg-category}；Keller, \emph{On differential graded categories}, arXiv:math/0601185。}
\end{conceptbox}

\begin{yizhu}[符号警告：$\bullet$ 与复合方向]
原文用缩小版的 bullet 表示复合运算 $\comp$，而公式里的 $\bullet$ 是分次指标（例如 $\hom^{\bullet}$、$\bigoplus_{\bullet\in\mathbb{Z}}$），两者外形相近但含义无关，不要混为一谈。更要留意复合方向：按上面的定义 $f\comp f'=(-1)^{-\vert f\vert}\mu_{\mathcal{A}}^2 (f',f)$，其中 $f\in\hom(L_0,L_1)$、$f'\in\hom(L_1,L_2)$，所以 $\comp$ 是「先 $f$ 后 $f'$」的图解顺序，与通常函数复合 $f'\circ f$ 的书写顺序恰好相反。后文所有 $\comp$ 表达式都按这一约定解读。
\end{yizhu}

\begin{dfn}
两个 $A_{\infty}$-范畴之间的一个\emph{$A_{\infty}$-函子} $\phi:\mathcal{A}\rightarrow\mathcal{B}$ 是一族映射 \[
    \left(\phi^d\colon\hom_{\mathcal{A}} (L_{d-1},L_d)\otimes\dotsb\otimes\hom_{\mathcal{A}} (L_0,L_1)\to\hom_{\mathcal{A}}^{\bullet +1-d} (\phi(L_0),\phi(L_d))\right)_{d\geq 1}
\] 连同对象层面的映射 $\phi\colon\Ob (\mathcal{A})\to\Ob (\mathcal{B})$，使得相容性方程
\begin{align*}
    \sum_r&\sum_{s_1+\dotsb+s_r=d}\mu_{\mathcal{B}}^r (\phi^{s_r}(a_d,\dotsb,a_{d-s_r+1}),\dotsb,\phi^{s_1} (a_{s_1},\dotsb, a_1))\\
    &=\sum_{m,n}(-1)^{\maltese_n}\phi^{d-m+1}(a_d,\dotsb, a_{n+m+1},\mu_{\mathcal{A}}^m (a_{n+m},\dotsb, a_{n+1}),a_n,\dotsb, a_1)
\end{align*}
成立。当 $\mathcal{A}$ 与 $\mathcal{B}$ 均为 dg 范畴时，若 $\phi^d\equiv 0$（$d\geq 2$），则称 $(\phi^d)$ 为\emph{dg 函子}。换言之，dg 函子可以看作对任意对象 $L,L'\in\Ob(\mathcal{A})$ 给出的上链映射 $\phi:\hom_{\mathcal{A}}(L,L')\rightarrow\hom_{\mathcal{B}}(\phi(L),\phi(L'))$，它保持复合（即 $\phi(f\comp f')=\phi(f)\comp \phi(f')$）与单位（即 $\phi(e_L)=e_{\phi(L)}$）。$A_{\infty}$/dg 函子的复合可按显然的方式定义。
\end{dfn}

\begin{yizhu}[原文笔误：$A_{\infty}$-函子的目标空间]
上面 $A_{\infty}$-函子定义中，映射 $\phi^d$ 的值域写作 $\hom_{\mathcal{A}}^{\bullet +1-d} (\phi(L_0),\phi(L_d))$，其下标应为 $\mathcal{B}$ 而不是 $\mathcal{A}$：$\phi^d$ 的取值落在目标范畴 $\mathcal{B}$ 的态射空间里，且 $\phi(L_0),\phi(L_d)$ 也说明这一点。另外，$\bullet+1-d$ 与后文写法 $[1-d]$ 是同一件事，都表示次数上移 $1-d$，不要误当成两种不同的移位约定。
\end{yizhu}

\begin{dfn}
对 $A_{\infty}$-范畴 $\mathcal{A}$，其\emph{上同调范畴} $H^{\bullet}(\mathcal{A})$ 是这样一个范畴：对象与 $\mathcal{A}$ 相同，而两个对象 $L_0, L_1$ 之间的态射空间为
$$
\Hom_{H^{\bullet}(\mathcal{A})}(L,L')
:=
H^{\bullet}(\hom_{\mathcal{A}}(L,L'),\mu_{\mathcal{A}}^1).
$$
$A_{\infty}$-函子 $(\phi^{\bullet}):\mathcal{A}\rightarrow\mathcal{B}$ 诱导普通函子 $H^{\bullet}(\phi):H^{\bullet}(\mathcal{A})\rightarrow H^{\bullet}(\mathcal{B})$，它在对象上的作用与 $\phi$ 相同，在态射上由 $[f]\mapsto [\phi^1(f)]$ 给出。

若诱导函子 $H^{\bullet}(\phi):H^{\bullet}(\mathcal{A})\rightarrow H^{\bullet}(\mathcal{B})$ 是范畴的嵌入（相应地，等价、同构），则称 dg 函子 $\phi:\mathcal{A}\rightarrow\mathcal{B}$ 为\emph{拟嵌入}（相应地，\emph{拟等价}、\emph{拟同构}）。
\end{dfn}
作为例子，我们描述一个由 $\field$-线性范畴构造出的特殊 dg 范畴。
\begin{dfn}
设 $\mathcal{C}$ 为 $\field$-线性范畴。$\mathcal{C}$ 中上链复形的\emph{dg 范畴 $\Ch(\mathcal{C})_{\dg}$} 定义如下：
\begin{itemize}
\item
\emph{对象}由 $\mathcal{C}$ 中的上链复形给出
$$
M
=
\left(\left(M^{\bullet}\right)_{\bullet\in\mathbb{Z}},
d_M:=\left(d_M^\bullet:M^{\bullet}\rightarrow M^{\bullet+1}\right)_{\bullet\in\mathbb{Z}}\right),
$$
其中 $M^{\bullet}$ 是 $\mathcal{C}$ 中的对象，$d_M^{\bullet}\in\Hom_{\mathcal{C}}(M^{\bullet},M^{\bullet+1})$ 是 $\mathcal{C}$ 中的态射，并且对每个 $\bullet\in\mathbb{Z}$ 满足 $d_M^{\bullet}\comp d_M^{\bullet+1}=0$。

\item
\emph{dg 态射空间}是 $\field$-模的上链复形，其第 $k$ 次分量为
$$
\hom_{\Ch(\mathcal{C})_{\dg}}^k
\left(
M,M'
\right)
:=
\prod_{\bullet\in\mathbb{Z}}
\Hom_{\mathcal{C}}\left(M^{\bullet},M'^{\bullet+k}\right),
$$
而微分为
$$
df:=d_M\comp f - (-1)^k f\comp d_{M'}
$$
其中 $f\in \hom_{\Ch(\mathcal{C})_{\dg}}^k(M,M')$。更具体地，$df$ 的第 $\bullet$ 次分量为
$$
(df)^{\bullet}
:=
d_M^{\bullet}\circ f^{\bullet+1}
-
(-1)^k
f^{\bullet}\circ d_{M'}^{\bullet+k}.
$$
\item
复合与单位按显然的方式定义。
\end{itemize}
我们记 $\Ch^b(\mathcal{C})_{\dg}$ 为 $\mathcal{C}$ 中\emph{有界上链复形的 dg 范畴}，其对象是 $\mathcal{C}$ 中满足除有限多个 $\bullet$ 外 $M^{\bullet}=0$ 的上链复形 $M=\left(\left(M^{\bullet}\right)_{\bullet\in\mathbb{Z}},d_M\right)$。态射空间的定义与 $\Ch(\mathcal{C})_{\dg}$ 中相同。
\end{dfn}
\begin{rmk}\label{rmk:EmbeddingOfCochainComplexes}
$\field$-线性范畴之间的函子 $\mathcal{F}:\mathcal{C}\hookrightarrow\mathcal{D}$ 按显然的方式自然地诱导 dg 范畴之间的 dg 函子
$
\mathcal{F}_{\dg}:
\Ch(\mathcal{C})_{\dg}
\hookrightarrow
\Ch(\mathcal{D})_{\dg}
$
与
$
\Ch^b(\mathcal{C})_{\dg}
\hookrightarrow
\Ch^b(\mathcal{D})_{\dg}.
$
若 $\mathcal{F}$ 是嵌入（相应地，等价、同构），则 $\mathcal{F}_{\dg}$ 是拟嵌入（相应地，拟等价、拟同构）。
\end{rmk}
\subsection[模与完美模]{$A_{\infty}$-模与完美模}
对 $A_{\infty}$-范畴 $\mathcal{A}$，回忆 \emph{$\mathcal{A}$-模} 是指 $A_{\infty}$-函子 $\mathcal{M}\colon\mathcal{A}\to\Ch_{\field}$，即从 $\mathcal{A}$ 到 $\field$ 上链复形的 dg 范畴的函子。
\begin{dfn}\label{dfn:A-inf-module}
     $\mathcal{M}\in\Fun (\mathcal{A},\Ch_{\field})\eqqcolon\Mod (\mathcal{A})$ 上的 $A_{\infty}$-模结构是一族 $\field$-线性映射
    \begin{align*}
    \mu^{d\vert 1}&\colon\hom_{\mathcal{A}} (L_{d-1},L_d)\otimes\dotsb\otimes\hom_{\mathcal{A}} (L_0, L_1)\otimes\mathcal{M} (L_0)\to\mathcal{M} (L_d) [1-d]
    \end{align*}
    使得对每个 $d\geq 0$，
% [原文 p.17]
\begin{equation}\label{eq:A-inf-bimodule-equation}\begin{aligned}
        &\sum (-1)^{\maltese_{d-i}}\mu_{\mathcal{M}}^{d-i\vert 1}(v_d,\dotsb,v_{i+1},\mu_{\mathcal{M}}^{i\vert 1} (v_i,\dotsb, v_1,\mathbf{b}))\\
        &+\sum (-1)^{\maltese_{d-k-i}}\mu_{\mathcal{M}}^{d-i+1\vert 1}(v_d,\dotsb, v_{k+i+1},\mu_{\mathcal{A}}^i (v_{k+i},\dotsb,v_{k+1}),v_k,\dotsb, v_1,\mathbf{b})=0.
    \end{aligned}\end{equation}
    设 $\mathcal{M}$ 与 $\mathcal{N}$ 为两个 $\mathcal{A}$-模。次数为 $k$ 的\emph{预态射} $s\colon\mathcal{M}\to\mathcal{N}$ 是一族与元组 $(L_0,\dotsb, L_{d-1})$ 相伴的链复形态射
    $s^d\colon\hom_{\mathcal A}(L_{d-2},L_{d-1})\otimes\dotsb\otimes\hom_{\mathcal A}(L_0,L_1)\otimes\mathcal M(L_0)\longrightarrow\mathcal N(L_{d-1})[k-d+1]$。

    从 $\mathcal{M}$ 到 $\mathcal{N}$ 的预态射全体构成 $\Bbbk$ 上的链复形，而 $\mathcal{A}$-模范畴 $\Mod (\mathcal{A})$ 是 dg 范畴，其微分在 \cite[page 20]{Seidel08} 中显式给出。满足 $\mu^1_{\Mod_{\mathcal{A}}} (s)=0$ 的预态射 $s$ 称为 $A_{\infty}$-模的\emph{态射}。若存在拟同构 $f\colon\mathcal{M}\to\mathcal{N}$，即存在态射 $g\colon\mathcal{N}\to\mathcal{M}$ 使得 $f\circ g$ 与 $g\circ f$ 取上同调后均为恒等，则称两个 $\mathcal{A}$-模 $\mathcal{M}$ 与 $\mathcal{N}$ \emph{同构}。
\end{dfn}
对每个对象 $X\in\Obj (\mathcal{A})$，对应 $X\mapsto\hom_{\mathcal{A}} (X,-)$ 按 \cite[section (1l)]{Seidel08} 中描述的方式定义一个 $A_{\infty}$-函子 $\mathcal{A}^{op}\to\Mod_{\field} (\mathcal{A})$。这是\emph{Yoneda 嵌入}的 $A_{\infty}$ 版本，记作 $\yo_{-}\colon\mathcal{A}\to\Mod (\mathcal{A})$。

\begin{yizhu}[为什么 Yoneda 嵌入送的是 $\mathcal{A}^{op}$]
固定对象 $X$ 后，$\hom_{\mathcal{A}}(X,-)$ 关于第二个变量是协变的；但把 $X$ 也看成变量时，$X$ 占据 Hom 的第一个槽，所以 $X\mapsto\hom_{\mathcal{A}}(X,-)$ 整体是反变的，给出的是 $\mathcal{A}^{op}\to\Mod(\mathcal{A})$ 的 $A_{\infty}$-函子，而非 $\mathcal{A}\to\Mod(\mathcal{A})$。这正是下面的完美模被说成 $\mathcal{A}^{op}$ 的幂等完备化的原因——$\Perf\mathcal{A}$ 作为 $\Mod(\mathcal{A})$ 的子范畴，从 $\mathcal{A}^{op}$ 的像出发生成。
\end{yizhu}

\begin{dfn}
    若存在对象 $X\in\Obj (\mathcal{A})$ 与 $\mathcal{A}$-模的同构 $\mathcal{M}\xrightarrow{\simeq}\yo_X$，则称 $A_{\infty}$-模 $\mathcal{M}\in\Mod_{\field} (\mathcal{A})$ 在 $\mathcal{A}$ 中\emph{可表}。

    我们定义\emph{完美 $\mathcal{A}$-模}的范畴 $\Perf\mathcal{A}\subseteq\Mod (\mathcal{A})$ 为包含 Yoneda 函子像的最小的预三角 dg 子范畴；等价地，$\Perf\mathcal{A}$ 是 $\mathcal{A}^{op}$ 的幂等完备化。
\end{dfn}

\begin{conceptbox}{完美模}
在 dg 范畴 $\mathcal{A}$ 上，模的导出范畴把 $\mathcal{A}$-模复形当作对象、把拟同构当作可逆态射；完美模 $\Perf\mathcal{A}$ 是其中「有限生成」的部分：从 Yoneda 像出发，在预三角意义下取最小闭包（允许平移、映射锥与有限直和），再补齐直和项，故等价地说就是 $\mathcal{A}^{op}$ 的幂等完备化。几何上它对应向量丛的复形：对概形 $X$ 有 $\Perf(X)=\Perf(\QCoh(X))$，由有限秩局部自由层生成。本文 A 侧出现的对象正是 $\Perf\mathcal{W}(X)^{op}$，它在 HMS 下与 B 侧的 $\Coh(Y)_{dg}$ 对应，因此主定理都陈述为 $\Perf\mathcal{W}(\cdots)$ 与 $\Coh(\cdots)_{dg}$ 之间的等价。
\par\smallskip\noindent{\footnotesize 参考：nLab，\emph{perfect complex}；Keller, \emph{On differential graded categories}, arXiv:math/0601185。}
\end{conceptbox}

